% page 145 of the facsimile (image p-144 of the scan)
% block 167.6 x 242.0 mm ; left margin 34.4 mm
\pgc{12.700mm}{0.000mm}{
\l{\cel{54}135}
\l{}
\l{}
\l{\sou{endogena}. (\sou{bot.}) Qua kreskas interne di su ipsa.}
\l{}
\l{\sou{endokardio}. (\sou{anat.}) Membrano qua tapetatre kovras la internajo,}
\l{\cel{3}la kavaji di la kordio.}
\l{}
\l{\sou{endokarpo}. (\sou{bot.)}\cel{2}Un ek la tri strati o membrani ye qui konsis-}
\l{\cel{3}tas la envelopilo (perikarpo) di la frukto, olta qua konstitu-}
\l{\cel{3}cas la surfaco interna.}
\l{}
\l{\sou{endoparazito}. Parazito qua habitas interne di sua hosto.}
\l{}
\l{\sou{endoplasmo}. Parto centrala od.interna di la korpo celula di la}
\l{\cel{3}enti monoplastida.}
\l{}
\l{\sou{endospermo}. Mikra organo, funcionanta kom kotiledono en la gra-}
\l{\cel{3}minei, la liliacei, e c.}
\l{}
\l{\sou{eneagono}. (\sou{geom.}) Poligono qua havas non edri e non anguli.}
\l{}
\l{\sou{enemiko}. Homo qua havas sentimenti des-amikala. - EFIS.}
\l{}
\l{\sou{energetiko}. Teorio fizikala, defensita da Duhem e Wilhelm}
\l{\cel{3}Ostwald, qua konsideras, en la fenomeni, nur la manifesti}
\l{\cel{3}mezurebla, abstraktante la naturo propra di la kozi. - DEFIRS.}
\l{}
\l{\sou{energio}. I. Forco vivaca, agera, di la organismo o di la psiko.}
\l{\cel{3}- II. Lo apta efektigar laboro. - DEFIRS.}
\l{}
\l{\sou{enfilar}. (\sou{trans.}) (\sou{Milito}). Kanonagar enemiki trupa o nava,}
\l{\cel{3}tale ke li trairesas longesale. - DEFIS.}
\l{}
\l{\sou{engajar}. (\sou{trans}.)Ligar ulu kom employoto per kontrato. - DEFIRS.}
\l{}
\l{\sou{engrosa}. (\sou{pri vendo o kompro di vari}). Granda-quantope.}
\l{}
\l{\sou{enigmato}.\cel{2}Ulo divinenda segun defino-frazi obskura.- EFIS.}
\l{}
\l{\sou{enkaustiko}.- I. (\sou{che la antiqui}). Farbo facita ek vaxi koloroza,}
\l{\cel{3}liquidigita per fairo. - II. Mixuro de vaxo e de terpentino,}
\l{\cel{3}quan onu uzas por briligar la parqueti, la mobli ligna, e c.}
\l{\cel{3}- III. Induturo vaxo-fundamenta, ye qua onu kovras la statui}
\l{\cel{3}marmora por igar plu matida lia blankeso. - DEFIS.}
\l{}
\l{\sou{enkilemo}. (\sou{citologio}).Liquido o suko nukleala en qua balnas}
\l{\cel{3}la grani kromatika.}
\l{}
\l{\sou{enklitiko}. (\sou{gram.}) Vorto qua, en la skribo, unionesas a la vorto}
\l{\cel{3}preiranta tale ke, aspekte, lu esas quaze una kun ita, e qua ne}
\l{\cel{3}acentizesas. Kom ex.: L.\cel{1}\sou{que}\cel{1}en\cel{1}\sou{neque}; ne en\cel{1}\sou{venisne}\cel{1}; F.}
\l{\cel{3}ce en\cel{1}\sou{est-ce}\cel{1}; je en sais-je. - DEFIS.}
\l{}
\l{enluminar. (\sou{trans}.)Piktar ye kolori vivaca, aplikata plata-nuance.}
\l{\cel{3}- DEFIS.}
}
